\section{Convex Hull Trick y Li Chao Tree}
\whenuse{Una DP o consulta debe calcular repetidamente
$\min_j(m_jx+b_j)$. Usar CHT con deque si pendientes y consultas son
mon\'otonas; Li Chao si llegan en cualquier orden.}

Cada estado anterior $j$ aporta una recta $y=m_jx+b_j$ y cada nuevo estado
consulta el menor valor en $x_i$. Esto reemplaza probar todos los $j$, que suele
ser $O(n^2)$. La plantilla incluida es Li Chao para m\'inimos: acepta rectas y
consultas desordenadas sobre un dominio entero conocido $[L,R]$.

\complexity{CHT mon\'otono: inserci\'on amortizada $O(1)$ y consulta $O(1)$;
Li Chao: $O(\log X)$ por operaci\'on y $O(q\log X)$ memoria en versi\'on din\'amica.}

\begin{lstlisting}[style=cpp]
struct LiChaoTree {
    struct Line {
        ll slope = 0, intercept = (1LL << 62);
        ll evaluate(ll x) const { return slope * x + intercept; }
    };
    struct Node { Line line; int left_child = -1, right_child = -1; };
    ll domain_left, domain_right;
    vector<Node> nodes = vector<Node>(1);

    LiChaoTree(ll left_bound, ll right_bound)
        : domain_left(left_bound), domain_right(right_bound) {
        assert(domain_left <= domain_right);
    }

    int make_node() {
        nodes.push_back(Node{});
        return (int)nodes.size() - 1;
    }
    void add_line(Line added) {
        add_line(0, domain_left, domain_right, added);
    }
    void add_line(int node, ll left, ll right, Line added) {
        ll middle = left + (ll)(((i128)right - left) / 2);
        bool better_left =
            added.evaluate(left) < nodes[node].line.evaluate(left);
        bool better_middle =
            added.evaluate(middle) < nodes[node].line.evaluate(middle);
        if (better_middle) swap(added, nodes[node].line);
        if (left == right) return;
        if (better_left != better_middle) {
            if (nodes[node].left_child == -1)
                nodes[node].left_child = make_node();
            add_line(nodes[node].left_child, left, middle, added);
        } else {
            if (nodes[node].right_child == -1)
                nodes[node].right_child = make_node();
            add_line(nodes[node].right_child, middle + 1, right, added);
        }
    }
    ll query(ll x) const {
        assert(domain_left <= x && x <= domain_right);
        return query(0, domain_left, domain_right, x);
    }
    ll query(int node, ll left, ll right, ll x) const {
        ll answer = nodes[node].line.evaluate(x);
        if (left == right) return answer;
        ll middle = left + (ll)(((i128)right - left) / 2);
        int next = x <= middle ? nodes[node].left_child
                               : nodes[node].right_child;
        if (next == -1) return answer;
        if (x <= middle) return min(answer, query(next, left, middle, x));
        return min(answer, query(next, middle + 1, right, x));
    }
};
\end{lstlisting}

\subsection{Ejemplo de uso}

Se agregan $y=2x+3$, $y=-x+10$ e $y=5$ y se consulta el m\'inimo.

\begin{lstlisting}[style=cpp]
void li_chao_example() {
    LiChaoTree lines(0, 10);
    lines.add_line({2, 3});   // 2x+3
    lines.add_line({-1, 10}); // -x+10
    lines.add_line({0, 5});   // 5

    cout << lines.query(2) << '\n'; // min(7,8,5)  = 5
    cout << lines.query(5) << '\n'; // min(13,5,5) = 5
    cout << lines.query(8) << '\n'; // min(19,2,5) = 2
}
\end{lstlisting}

\warning{Puede desbordar $mx+b$; usar \texttt{\_\_int128} si los l\'imites lo
requieren y elegir un infinito mayor que toda respuesta v\'alida. Para m\'aximos,
negar rectas/respuestas o invertir las comparaciones.}
